\abstract[Abstract]{%
The ICH E17 guideline frames regional pooling in multi-regional clinical trials (MRCTs) as a question of effect modifier distributional similarity, but supplies no methodology. We propose a framework measuring dissimilarity with the Wasserstein-1 ($W_1$) distance, retained in each effect modifier's original units. Unlike the standardized mean difference (SMD), $W_1$ responds to differences in scale and shape as well as location, and needs only baseline distributions, so it can be computed at planning from prior trials, registries, or real-world evidence. Calibration follows two pathways: given prior evidence on conditional average treatment effect (CATE) sensitivity ($L$), $W_1$ becomes a maximum potential regional treatment effect difference ($\Delta_{\max}$) with bootstrap confidence intervals; absent such evidence, the required sensitivity ($L^*$) is reverse-calculated at pre-specified clinical margins. A decision study compared six measures across four simulated worlds: every summary-based competitor was at chance through 2000 patients per region in at least one clinically plausible world, whereas $W_1$ recovered the true partner structure in all four, and in the mean-matched cells reached the accuracy target at a fraction of the Kolmogorov--Smirnov statistic's required sample size. An estimation study across seven scenarios showed satisfactory bias and coverage of the percentile bootstrap for non-negligible distributional differences at moderate per-region sample sizes. A hypothetical thrombolytic MRCT using GUSTO-I patient data ranked 15 partners against a small-sample anchor on age and systolic blood pressure; rankings diverged markedly between the two, showing that modifiers must be evaluated jointly. The framework converts ``similar enough'' into a reproducible basis for sponsor judgment.
}
